\documentclass[11pt,a4paper]{article}
\usepackage[T1]{fontenc}
\usepackage[utf8]{inputenc}
\usepackage{lmodern}
\usepackage{amsmath,amssymb,booktabs}
\usepackage[hidelinks]{hyperref}
\title{Sparse-Channel EEG Affect Representation Under Cross-Dataset and Montage Shift}
\author{}
\date{}
\begin{document}
\maketitle
\begin{abstract}
Does a low-channel model remain calibrated on held-out people and across corpora when reference montage and acquisition domains change? The manuscript defines a reproducible research design that tests an explicit methodological contribution against appropriate non-adaptive or low-capacity baselines. All comparisons are prospective; no empirical improvement or validated model performance is claimed at this stage.
\end{abstract}
\section{Introduction and related methodological problem}
Does a low-channel model remain calibrated on held-out people and across corpora when reference montage and acquisition domains change? The scientific goal is to distinguish a genuine computational contribution from gains due to subject overlap, dataset bias, confounded acquisition settings, or favourable synthetic conditions.
\section{Research objectives and hypotheses}
Channel sparsity  by  corpus shift factorial benchmark with explicit montage-reference compatibility, transfer-risk uncertainty and no false Muse-equivalence claims. Hypothesis H1 predicts better pre-specified robustness or accuracy than an appropriate baseline under independently held-out evaluation; hypothesis H2 requires that the measured gain remains interpretable after accounting for uncertainty and deployment constraints. Both hypotheses are proposals, not tested findings.
\section{Materials and data eligibility}
DEAP, DREAMER and AMIGOS (only after licence/data audit); open-licensed fallback candidates GAMEEMO v3 and FACED NEMAR v1.1.4 (also requiring actual file and reference QA). Public EEG only. No local volunteers necessary. Data usage requires source licensing, stable event/label definitions, and a documented eligibility gate. Any unavailable or incomparable source must be excluded with an explicit reason.
\section{Methods and experimental design}
Create montage registry; compare available channel subsets (not claiming Muse channel equivalence), spectral/entropy features and compact CNN; grouped LOSO in each corpus, external corpus only if labels/reference harmonize. Compare with full-channel, linear and shuffled-label baselines. Hyperparameters and calibration are selected only from training splits. Evaluation units and confidence intervals must correspond to independent subjects, sessions, devices or scenarios rather than artificially replicated samples.
\section{Statistical analysis and primary outcomes}
Macro-F1, balanced accuracy, AUROC where applicable, concordance for continuous targets, ECE, Brier, CPU time, clustered 95\% CI. Every primary comparison will report uncertainty intervals, failure cases, negative controls and the actual sample/split sizes. Planned sensitivity analyses vary missingness, corruption and acquisition conditions.
\section{Results}
No empirical dataset analysis or confirmatory comparison has been completed for this manuscript. This section will report results only after authorised data collection or reproducible benchmark execution; planned metrics are not observed effects.
\section{Discussion, limitations and continuation criteria}
Montage locations may differ; stimulus, sampling and label semantics confound transfer; restricted datasets may not be available. Show reproducible subject-disjoint evaluation, positive delta against appropriate low-channel baseline with CIs, and explicit limitations for any non-equivalent montage. A novel method cannot be asserted solely from a different application domain; comparison against current prior art and reproducibility on an independent setting are required.
\section{Reproducibility and provenance}
The original research plan is preserved in \texttt{source\_plan.md}; repository and exact source commit are recorded in \texttt{provenance.json}. No sensitive recordings or private credentials are embedded.
\end{document}
